\chapter{Boundary systems, dimensional range, and action functional}
\label{chap:pantalla}

In Dimensional Mechanics, dimension is obtained from the number of
independent channels available at a boundary. It is not appended as an external label
to the state space: it is computed as the rank of a module of prolongations
subject to its relations. A change in dimension then corresponds to a
variation in rank during transport.

\section{Constraint systems and boundary data}

Let
\[
 \cdots\longrightarrow\mathcal S_{n+1}
 \xrightarrow{r_{n+1,n}}\mathcal S_n\longrightarrow\cdots
\]
an inverse system of surviving states. The boundary structure at
level \(n\) is the pair
\[
 \mathcal F_n=(r_{n+1,n},\beta_n),
\]
where \(\beta_n\) preserves the minimum datum necessary to decide the
compatibility of a prolongation. Its coordinates may include the arithmetic
type, the orientation, the carry cocycle, and the incidence class.

\begin{lemma}[Compatibility of successive constraints]
Suppose
\[
 r_{n+2,n}=r_{n+1,n}\circ r_{n+2,n+1},
 \qquad
 \beta_n\circ r_{n+1,n}
 =b_{n+1,n}\circ\beta_{n+1}.
\]
Then every compatibility condition expressed through \(\beta_{n+1}\)
induces a compatible condition at level \(n\).
\end{lemma}
\begin{proof}
For \(x\in\mathcal S_{n+2}\), equality of compositions identifies the
direct restriction with the two consecutive restrictions. The second
equality transports the boundary datum through the same diagram. The
condition at level \(n\) is therefore the image of the deeper
condition.
\end{proof}

\section{Module of channels and dimensional range}

Fixing a field \(\Kfield_n\), let \(M_n(x)\) be the module generated by the
prolongation channels of \(x\), and let \(\mathcal R_n(x)\) be the submodule of
relations among them.

\begin{definition}[Dimensional screen]
The screen associated with \(x\) at level \(n\) is
\[
 \Sigma_n(x)=
 (\Kfield_n,M_n(x),\beta_n(x),\mathcal R_n(x)),
\]
and its effective dimension is
\[
 d_{\Sigma_n}(x)=
 \dim_{\Kfield_n}\frac{M_n(x)}{\mathcal R_n(x)}.
\]
If \(M_n(x)\cong\Kfield_n^{g_n}\) and the columnas of a matrix
\(R_n(x)\) generate \(\mathcal R_n(x)\), then
\[
 d_{\Sigma_n}(x)=g_n-\rank_{\Kfield_n}R_n(x).
\]
When a complement is fixed and \(P_n(x)\) denotes the projector whose image is
isomorphic to the quotient, the same dimension is
\(d_{\Sigma_n}(x)=\rank_{\Kfield_n}P_n(x)\).
\end{definition}

\begin{proposition}[Invariance under equivalence of presentations]
The screen dimension is invariant under invertible changes of basis in the generators and in the relations.
\end{proposition}
\begin{proof}
A change of bases transforms the relation matrix \(R\) into \(URV\), with
\(U\) and \(V\) invertible. Since \(\rank(URV)=\rank R\), the number
\(g-\rank R\) does not change. Equivalently, two projectors presenting the
same quotient have isomorphic images and therefore the same rank.
\end{proof}

For a transport \(T_\gamma:\mathcal S_n\to\mathcal S_m\), define
\[
 \Delta d(\gamma;x)=
 d_{\Sigma_m}(T_\gamma x)-d_{\Sigma_n}(x).
\]
A dimensional transition simultaneously transports the state and its channel
module; fixed-rank dynamics is the case \(\Delta d=0\).

\section{Inherited arithmetic and lattice screens of HMT}

The ranks that occur below are not labels appended to the
equations. They arise from distinct linear constructions and preserve their
selection maps. The punctured ternary Golay code
\[
 G_{11}=[11,6,5]_3\subset\mathbb F_3^{11}
\]
is perfect of radius two \cite{MacWilliamsSloane1977}, since the balls
centered at its \(3^6\) words are
disjoint and
\[
 V_3(11,2)=1+2\binom{11}{1}+2^2\binom{11}{2}
 =1+22+220=243,
\qquad
 3^6\cdot243=3^{11}.
\]
Here \(243\) is the cardinal of the syndrome space. It is not identified by
mere numerical equality with the layer of \(243\) cells of the cubic
completion of APP\@.

The second chain starts from the permutation module
\(V_{12}=\mathbb R^{12}\) of \(A_5\) on the twelve vertices of the icosahedron.
To fix the gauge, we let \(A_5\) act on
\(\{0,1,2,3,4\}\), take
\(H=\langle(0\ 1\ 2\ 3\ 4)\rangle\cong C_5\) and identify the position
\(j\) with the left coset \(r_jH\). In image notation, the twelve
representatives are
\[
\begin{array}{c|c@{\qquad}c|c@{\qquad}c|c}
j&r_j&j&r_j&j&r_j\\ \hline
1&(3,2,4,1,0)&5&(3,2,0,4,1)&9&(3,1,2,4,0)\\
2&(3,1,0,2,4)&6&(4,0,1,2,3)&10&(2,1,4,3,0)\\
3&(0,1,3,4,2)&7&(4,3,2,1,0)&11&(3,4,1,2,0)\\
4&(1,4,2,3,0)&8&(0,2,3,1,4)&12&(4,2,1,3,0).
\end{array}
\]
We denote by \(5A\) the conjugacy class containing the generator
\(g_5=(0\ 1\ 2\ 3\ 4)\) of \(H\), and by \(5B\) the class containing
\(g_5^2\). In the class order \(1A,2A,3A,5A,5B\), the permutation character is
\(\chi_{12}=(12,0,0,2,2)\). For
\(\phi=(1+\sqrt5)/2\) and \(\phi'=(1-\sqrt5)/2\), we choose the irreducible
character
\[
 \chi_3=(3,-1,0,\phi,\phi').
\]
The scalar product of characters gives the exact decomposition.
\[
 V_{12}\cong V_1\oplus V_3\oplus V_{3'}\oplus V_5.
\]

The projector
\[
 P_{11}=I_{12}-\frac1{12}J_{12}
\]
eliminates exclusively the uniform vector and has image
\[
 V_{11}=\mathbf1^\perp\cong V_3\oplus V_{3'}\oplus V_5.
\]
For the dodecaphasic state
\[
 K=(234,543,140,729,659,824,621,58,914,794,146,601),
 \qquad \sum_jK_j=6263,
\]
let \(\widetilde K=P_{11}K\). The preceding position--coset correspondence
defines the permutation matrices \(\rho(g)\), and the central projector
\[
 P_3=\frac3{60}\sum_{g\in A_5}\chi_3(g^{-1})\rho(g)
\]
satisfies \(P_3^2=P_3=P_3^{\mathsf T}\) and has rank three. Exact evaluation
at \(\mathbb Q(\sqrt5)\) produces
\[
 u_K=P_3\widetilde K,
 \qquad
 \|u_K\|^2=\frac{6638585+2275584\sqrt5}{20}>0.
\]
Consequently, the line \(\ell_K=\mathbb Ru_K\) is well defined and
\[
 V_{10}:=(V_3\cap\ell_K^\perp)\oplus V_{3'}\oplus V_5,
 \qquad
 V_{11}=\ell_K\oplus V_{10},
 \qquad \dim V_{10}=10.
\]
The chain \(12\to11\to10\) means, respectively, removing the
uniform component and the line polarized by \(K\). The finite set of
sixty elements is divided into twelve classes; its enumeration
reconstructs the projector and norm exactly.

The reduction \(26\to24\) belongs to another category. Let \(U\) be the integral
hyperbolic plane with Gram matrix
\(\left(\begin{smallmatrix}0&1\\1&0\end{smallmatrix}\right)\), y sea
\[
 L_{26}=\Lambda_{24}\oplus U.
\]
With the hyperbolic summand fixed, this is the standard presentation of the
even unimodular Lorentzian lattice \(II_{25,1}\)
\cite{ConwaySloane1999}.
For a primitive isotropic vector \(e_+\in U\),
\[
 e_+^\perp=\Lambda_{24}\oplus\mathbb Ze_+,
\qquad
 e_+^\perp/\mathbb Ze_+\cong\Lambda_{24}.
\]
The decrease of rank by two comes from imposing orthogonality and then quotienting
by the null line. It is not the reduction \(11\to10\) repeated with other
numbers.

\section{Involution between residual data and Euclidean quotient}

The decomposition \(n=9Q+r\) suggests two coordinates of different kinds:
the residual representative \(r\) and the quotient \(Q\). To formulate their
exchange on a stable domain, we consider
\[
 \mathcal D=\mathbb Z^2\times\mathbb R_{>0},
\qquad
 E_R(m,w)=\frac{m^2}{R^2}+w^2R^2,
\]
and we define
\[
 \mathsf T(m,w,R)=(w,m,R^{-1}).
\]
\begin{proposition}[Residual--Euclidean Duality]
\(\mathsf T\) is an involution and
\[
 E_R(m,w)=E_{R^{-1}}(w,m).
\]
On \(\ell^2(\mathbb Z^2)\), the unitary
\(U_T|m,w\rangle=|w,m\rangle\) satisfies
\[
 U_TH(R)U_T^{-1}=H(R^{-1}),
\qquad
 H(R)|m,w\rangle=E_R(m,w)|m,w\rangle.
\]
\end{proposition}
\begin{proof}
The second map again exchanges the coordinates and inverts
\(R\) twice. The energy identity is obtained by direct substitution; the
operatorial identity is checked on the orthonormal basis.
\end{proof}

The proposition is an exact lattice duality. The physical reading as
momentum, winding, and compactification radius requires an additional map
to the state space of a string theory; it does not follow from the
algebraic involution \cite{Giveon1994}.

\section{Typed interface with undecadimensional reduction}

The screen \(V_{11}=\ell_K\oplus V_{10}\) and the involution
\((m,w,R)\mapsto(w,m,R^{-1})\) provide the two algebraic data that
motivate comparison with the M-theory context: a marked reduction
from eleven to ten coordinates and an equivalence of reciprocal radii. For
that comparison to be a physical realization, a morphism must be specified
\[
 \mathcal R_M:
 (\mathcal X_{\rm TPK},V_{11},\ell_K,\mathsf T)
 \longrightarrow
 (\mathcal H_M,\mathcal F_{10},R_{\rm phys},\alpha',
 \mathcal Q_{\rm susy}),
\]
where the codomain includes a state space, ten-dimensional
fields, physical radius, string tension, and supersymmetric charges
\cite{Witten1995}. The algebraic decomposition fixes the domain of
\(\mathcal R_M\); it does not by itself introduce oscillators, branes,
supersymmetry, or the supergravity equations.

This separation preserves the depth of the correspondence without altering
the types. The reduction \(11\to10\), the involution of radii, and the quotient
\(26\to24\) are theorems of the objects defined. The assertion that a
concrete dynamical theory is their image requires exhibiting \(\mathcal R_M\) and proving
that it intertwines actions, restrictions, and observables.

\section{Dimensional state and simultaneous transport}

The state on which Dimensional Mechanics acts is written
\[
 \mathfrak x=(x,\Sigma_n(x),\beta_n(x),[\gamma]),
\]
where \(x\in\mathcal X_{\rm TPK}\), \(\Sigma_n(x)\) is its screen,
\(\beta_n(x)\) is the boundary datum, and \([\gamma]\) is the class of paths
compatible with that datum. A transport does not act only on \(x\): it consists of
\[
 (T_\gamma,F_\gamma):
 (x,\Sigma_n,\beta_n)\longrightarrow
 (T_\gamma x,\Sigma_m,\beta_m),
\]
where \(F_\gamma:M_n(x)\to M_m(T_\gamma x)\) maps relations to relations
and satisfies the compatibility diagram with \(\beta_n,\beta_m\). The dimensional
jump is the change in rank of the induced morphism
\[
 \overline F_\gamma:
 M_n(x)/\mathcal R_n(x)\longrightarrow
 M_m(T_\gamma x)/\mathcal R_m(T_\gamma x).
\]
This formulation prevents assigning a dimension to the isolated point: the dimension
belongs to the system of channels that accompanies the state and its transport.

\section{Descent to a dynamics on point states}

Let \(q_{\rm pt}:\mathcal X_{\rm TPK}\to Z\) be a projection that eliminates the
genealogy of the state, and define
\[
 Y_{\rm pt}:=\operatorname{im}(q_{\rm pt}).
\]
Let, moreover, \(U:\mathcal X_{\rm TPK}\to\mathcal X_{\rm TPK}\) a
evolution complete.

\begin{theorem}[Fiberwise Descent Criterion]
There exists a unique map
\(\overline U:Y_{\rm pt}\to Y_{\rm pt}\) such that
\(q_{\rm pt}U=\overline Uq_{\rm pt}\) if and only if
\[
 q_{\rm pt}(x)=q_{\rm pt}(y)
 \Longrightarrow q_{\rm pt}(Ux)=q_{\rm pt}(Uy).
\]
\end{theorem}
\begin{proof}
Necessity results by applying \(\overline U\) to two representatives of one and the
same fiber. For sufficiency, one defines
\[
 \overline U(q_{\rm pt}(x))=q_{\rm pt}(Ux).
\]
The condition guarantees that the definition does not depend on the representative. Since
the codomain has been fixed equal to the image of \(q_{\rm pt}\), the
surjectivity required for uniqueness is automatic.
\end{proof}

When the condition fails, a single point class has distinct images;
the descended evolution is then multivalued. Retaining the memory datum,
restricting the domain, or working with a correspondence are three distinct
mathematical modifications.

\section{Additive functional on oriented trajectories}

To avoid a hidden path-composition term, the state of an edge is
extended by the direction and arithmetic type of the preceding edge. If
\(\widetilde e_j=(e_j,d_{j-1},t_{j-1})\), define the local weights
\[
 w_{\rm giro}(\widetilde e_j)=\mathbf1_{d_j\ne d_{j-1}},
 \qquad
 w_{\rm tipo}(\widetilde e_j)=\mathbf1_{t_j\ne t_{j-1}}.
\]
The data \((d_0,t_0)\) form part of the initial boundary condition.
For an augmented trajectory
\(\widetilde\gamma=(\widetilde e_1,\ldots,\widetilde e_R)\), let
\[
 N_{\rm giro}(\widetilde\gamma)=\sum_{j=1}^R
 w_{\rm giro}(\widetilde e_j),
 \qquad
 N_{\rm tipo}(\widetilde\gamma)=\sum_{j=1}^R
 w_{\rm tipo}(\widetilde e_j).
\]
We define
\[
 \mathbf I(\widetilde\gamma)=
 (N_{\rm giro}(\widetilde\gamma),N_{\rm tipo}(\widetilde\gamma)).
\]
The composition is allowed only when the preceding datum of the second path
coincides with the terminal datum of the first. With this convention, the partition
of the sum into the two edge intervals proves
\[
 \mathbf I(\widetilde\gamma\widetilde\eta)
 =\mathbf I(\widetilde\gamma)+\mathbf I(\widetilde\eta),
\]
so that \(\mathbf I\) is a cocycle with values in \(\NN^2\). For
\(\lambda>0\), the combination
\[
 I_\lambda(\widetilde\gamma)
 =N_{\rm giro}(\widetilde\gamma)
 +\lambda N_{\rm tipo}(\widetilde\gamma)
\]
is a scalar functional over the path space.

In the ordinary Lagrangian formulation, the action functional is defined
on curves in a previously selected configuration space. Here the
order is prior: \(I_\lambda\) quantifies changes in direction and arithmetic law
in the category of augmented trajectories itself. A realization map
may subsequently transport this cocycle into a continuous action. The
additivity proved above is the property that permits that
comparison without identifying the two functionals in advance.

\begin{principle}[Minimal interaction]
Among admissible trajectories with fixed endpoints and boundary data, the
selected evolution minimizes \(I_\lambda\).
\end{principle}

\begin{proposition}[Finite Existence of Minimizers]
If \(\Gamma(a,b)\) is finite and nonempty, \(I_\lambda\) attains a minimum in
\(\Gamma(a,b)\).
\end{proposition}
\begin{proof}
The image of a non-empty finite set under a real function has a minimum element.
\end{proof}

The mass character of \cref{chap:masas} constitutes another functional on
trajectories: it transforms an integer cocycle into a positive ratio. The
functional and the mechanism selecting the trajectory are distinct objects;
their composition will be specified before introducing a unit of
mass.
